\subsection{Arithmetic norm review and comparable-degree windows}\label{endpoint-pr25_norm_original--arithmetic-norm-review-and-comparable-degree-windows}

The owner supplied the 2026-09-13 arithmetic endpoint note. Its principal input is the actual density w\_h=\textbar2xi/h\textbar\^{}2/(2pi), its k-fold convolution, and the unscaled monic squared norms omega\_(h,k,j). This is a review and extension of that written proof, not a Lean certificate for zeta estimates or arithmetic integrals. The locally linked delivery archive was not accessible in this runtime; its reported test and manifest results are not counted as fresh executions.

\subsubsection{1. Audit of the analytic steps}\label{endpoint-pr25_norm_original--audit-of-the-analytic-steps}

The Joukowski ellipse with \textbar w\textbar=8 has semiaxes 65/16 and 63/16. For f\_T(z)=zeta(1/2+iT+iz), \textbar T\textbar\textgreater=10 keeps the pole outside its closure. Its image lies in the fixed real strip {[}-55/16,71/16{]}, with imaginary magnitude between \textbar T\textbar-65/16 and \textbar T\textbar+65/16. Split that strip at Re(s)=1/4. Above this value, the fractional-part integral gives polynomial growth O((1+\textbar T\textbar)); below it the functional equation and vertical-strip gamma estimate give O((1+\textbar T\textbar)\^{}(1/2-Re(s)+1)). The largest exponent is 79/16\textless6. Thus the claimed envelope M\_T=C\_0(1+\textbar T\textbar)\^{}6 is sufficient, uniformly on the ellipse. Increase C\_0 so M\_T\textgreater=1.

The point z=-3i/2 corresponds to \textbar w\textbar=(3+sqrt(13))/2\textless4 and hence to zeta(2+iT). The reciprocal Dirichlet series gives modulus at least zeta(2)\^{}(-1). The annular maximum principle gives m\_T\textgreater=zeta(2)\^{}(-3) M\_T\^{}(-2). Every closed unit disk centred on {[}-1,1{]} lies strictly within the ellipse: its surrounding rectangle {[}-2,2{]}x{[}-1,1{]} already satisfies the ellipse inequality. The maximum principle and Cauchy estimate therefore give sup\textbar f\_T'\textbar\textless=M\_T. At a maximizer use the longer one-sided interval in {[}-1,1{]}, of length at least one; m\_T/(2M\_T)\textless=1/2. Integration gives m\_T\^{}3/(8M\_T) and the stated exponent 42. The compact remaining T-range is handled by continuity and analytic nonvanishing on an interval. No nonvanishing on the critical line is assumed.

For the passage to \textbar g/h\textbar, use g=h(g/h) before taking absolute values. The polynomial upper bound on \textbar h\textbar{} avoids division at a selected zero. Gamma asymptotics yield the claimed interval-mass bound. One regularity detail should be explicit: w\_h is bounded as well as integrable, by its continuous entire quotient on compact sets and exponential tail. Thus w\_h*w\_h is continuous, using L1 translation continuity against the bounded factor. Its integrand is positive almost everywhere, so the convolution is strictly positive everywhere. This justifies its positive minimum on {[}-1,1{]}. Restricting the remaining factors to that interval proves the note's lower envelope with the original factor vartheta\_h\^{}(k-3).

The monic Legendre norm on {[}-L,L{]} is exactly

ENDPOINTSOURCE072CC5FCCE66B5C5698D8622SLOT0END

The trial upper bound at degree j is (2j)! b\^{}(-2j)(M\_h(b)\^{}k+M\_h(-b)\^{}k), for fixed 0\textless b\textless pi/2. Both are used on the original shifted line S=k/2+iu; multiplication by i\^{}j has modulus one and does not change the leading coefficient norm. These arguments validate the supplied O\_h(n) doubling-window result as a written theorem, not as a floating-point or Lean assertion.

\subsubsection{2. Stronger, two-sided envelope}\label{endpoint-pr25_norm_original--stronger-two-sided-envelope}

Write a=pi/2 and retain the lower-envelope constants c\textgreater0, vartheta\textgreater0, B=42+2 deg(h). Put

ENDPOINTSOURCE072CC5FCCE66B5C5698D8622SLOT1END

Then, for every n\textgreater=k\textgreater=3,

ENDPOINTSOURCE072CC5FCCE66B5C5698D8622SLOT2END

Proof of the lower bound: l\_n(n)\textgreater=(2/3)n\^{}(2n)4\^{}(-n); vartheta\^{}(k-3)\textgreater=t\^{}n; exp(-a(n+k-3))\textgreater=exp(-2an); (1+n+k-3)\^{}(-B)\textgreater=(2n)\^{}(-B)\textgreater=2\^{}(-nB), since 2n\textless=2\^{}n for integer n\textgreater=1; finally 2c/3\textgreater=c0\^{}n.~Multiply. For the upper bound use (2n)!\textless=(2n)\^{}(2n), k\textless=n and M\_h(b)\^{}k+M\_h(-b)\^{}k\textless=2M\^{}n, then 2\textless=4\^{}n.~This proves (A), retaining all masses in constants rather than resetting them.

For n\textgreater=k\textgreater=3 and n/2\textless=r\textless=2n, let m=n+r. From (A),

ENDPOINTSOURCE072CC5FCCE66B5C5698D8622SLOT3END

This extends the doubling-window estimate to the overlapping windows needed by the volume theorem. Constants depend only on the fixed h and chosen b, not k,n,r. No monotonicity of the monic norms is assumed. The inequality is valid at r=n and at r=n+1 when n\textgreater=1.

References: NIST DLMF 25.4 (functional equation), 5.11.9 (uniform bounded-real-part gamma asymptotic), 18.3 (Legendre normalization). The general arguments above are written proofs; the present Lean contribution formalizes their finite norm-to-volume consequences, not their analytic hypotheses.
